\documentclass[11pt,a4paper]{article}
\usepackage[T1]{fontenc}
\usepackage[{{babel}}]{babel}
\usepackage{lmodern}
\usepackage[margin=2.2cm, headheight=14pt]{geometry}
\usepackage{amsmath, amssymb, amsthm}
\usepackage{enumitem}
\usepackage{fancyhdr}
\usepackage[hidelinks]{hyperref}

% {{fr:En-tête et pied de page|en:Header and footer}}
\pagestyle{fancy}
\fancyhf{}
\lhead{{{author}}}
\chead{{{title}}}
\rhead{{{institution}}}
\cfoot{\thepage}

% {{fr:Un environnement par exercice|en:One environment per exercise}}
\newcounter{exercice}
\newenvironment{exercice}[1][]{%
  \refstepcounter{exercice}\par\medskip
  \noindent\textbf{{{fr:Exercice|en:Exercise}}~\theexercice}\if\relax\detokenize{#1}\relax\else~(#1)\fi\par\smallskip}{\par\medskip}

\newcommand{\R}{\mathbb{R}}
\newcommand{\N}{\mathbb{N}}

\begin{document}

\begin{center}
  {\LARGE\bfseries {{title}}}\\[4pt]
  {{author}} --- {{date}}
\end{center}

\begin{exercice}[{{fr:Suites|en:Sequences}}]
  \label{ex:suites}
  \begin{enumerate}[label=\alph*)]
    \item {{fr:Montrer que la suite $(u_n)_{n \in \N}$ définie par $u_n = \frac{1}{n+1}$ converge.|en:Show that the sequence $(u_n)_{n \in \N}$ defined by $u_n = \frac{1}{n+1}$ converges.}}

    \begin{proof}
      {{fr:Pour tout $\varepsilon > 0$, on choisit $N > \frac{1}{\varepsilon}$. Alors pour $n \geq N$, $|u_n| < \varepsilon$.|en:For every $\varepsilon > 0$, choose $N > \frac{1}{\varepsilon}$. Then for $n \geq N$, $|u_n| < \varepsilon$.}}
    \end{proof}

    \item {{fr:Calculer la limite.|en:Compute the limit.}}
    \[
      \lim_{n \to \infty} u_n = 0
    \]
  \end{enumerate}
\end{exercice}

\begin{exercice}
  {{fr:Résoudre le système suivant dans $\R^2$ (voir aussi l'exercice~\ref{ex:suites}).|en:Solve the following system in $\R^2$ (see also exercise~\ref{ex:suites}).}}
  \begin{align*}
    2x + y &= 3 \\
    x - y  &= 0
  \end{align*}
  {{fr:Solution :|en:Solution:}} $x = y = 1$.
\end{exercice}

\end{document}
